% ===========================================================================
%  Chapter 1: Introduction
%  Dual-Target RL Molecular Generation — Workshop Report
% ===========================================================================

Polypharmacology---the deliberate modulation of multiple biological
targets by a single molecule---has emerged as a promising strategy for
complex multifactorial diseases~\citep{bolognesi2013polypharmacology,
anighoro2014polypharmacology}.
In Alzheimer's disease alone, the simultaneous inhibition of
cholinesterases (AChE, BChE), $\beta$-secretase (BACE1), and kinases
involved in tau phosphorylation (GSK3$\beta$, CDK5) has been pursued
both computationally and
clinically~\citep{stern2022multitarget, weinreb2012cholinesterase}.
Approved dual-target agents such as lapatinib (EGFR/HER2) and
ladostigil (AChE/MAO-B) demonstrate that single molecules \emph{can}
engage structurally distinct binding sites, but their discovery has
relied on serendipity or labor-intensive medicinal chemistry campaigns.

Reinforcement learning (RL) offers a principled framework for
goal-directed molecular generation: an agent constructs molecules
step by step and receives a scalar reward derived from a computational
scoring function~\citep{telepov2023freedpp, loeffler2024reinvent4}.
To date, RL-based generators have been applied almost exclusively
in single-target mode, optimizing docking score or predicted affinity
toward one protein at a time.
Multi-objective extensions typically optimize \emph{property}
trade-offs (QED, SA, logP) rather than binding affinity to two
specific protein targets
simultaneously~\citep{fromer2023moo, wang2024scarlp}.

This study addresses three gaps identified in our systematic
literature survey (Section~\ref{sec:rw:moo}):
\begin{enumerate}
\item \textbf{Systematic weight-space exploration.}
  No prior work scans the continuous weight space from selective
  ($w_A{=}1, w_B{=}{-1}$) through single-target ($w_A{=}1, w_B{=}0$)
  to dual-target ($w_A{=}1, w_B{=}1$) using structure-based docking
  as the RL reward for two protein targets.
  We implement six reward configurations---\texttt{single\_a},
  \texttt{single\_b}, \texttt{double}, \texttt{linear\_equal},
  \texttt{biased\_a}, and \texttt{selective\_a}---within the same
  generator architecture (FREED++).
\item \textbf{Orthogonal evaluation.}
  Most generative studies evaluate molecules with the same scorer
  used as the RL reward, creating a circular evaluation.
  We decouple training (UniDock docking scores $d_A, d_B$) from
  evaluation (Boltz-2 affinity predictions $a_A, a_B$), using two
  fundamentally different methods---physics-based empirical scoring
  versus a transformer-based learned surrogate---to mitigate
  scorer-specific artifacts.
\item \textbf{Selectivity--polypharmacology spectrum.}
  The literature treats target selectivity and multi-target design
  as separate research fields.  We show that these are points on a
  single parametric scale controlled by the reward weight vector,
  enabling direct comparison within a unified experimental framework.
\end{enumerate}

The primary hypothesis is that dual-target RL generation
(\texttt{double} mode, $R = -\max(d_A, d_B)$) increases the
probability of obtaining simultaneous hits on both targets relative
to single-target baselines (\texttt{single\_a}, \texttt{single\_b}),
as measured by Boltz-2 affinity evaluated independently of the RL
reward.

The experiment is organized as a fully reproducible eight-stage
pipeline: data acquisition, target preparation, calibration,
generation, Boltz-2 scoring, aggregation and statistical
testing, report generation, and claim audit.
All comparisons enforce fairness through matched molecular budgets,
fixed seed sets, shared hit-threshold definitions, and uniform
post-processing rules.
Success criteria, pair assignments, and the statistical analysis
plan are pre-registered before any generation runs.
